%%%%%%%%%%%%%%%%%%%%%%%%%%%%%%%%%%%%%%%%%%%%%%%%%%%%%%%%%%%%%%%%%%%%%%
% Lớp 11 | Đạo hàm | Đạo hàm cấp hai
% Dạng 1D7?3-1: Tính đạo hàm cấp hai
% Số câu: 40  (sinh tự động bởi tach_cau.py; từ nay thầy sửa trực tiếp file này)
%%%%%%%%%%%%%%%%%%%%%%%%%%%%%%%%%%%%%%%%%%%%%%%%%%%%%%%%%%%%%%%%%%%%%%

%%%%%%%%%%%%%%%% TRẮC NGHIỆM 4 PHƯƠNG ÁN %%%%%%%%%%%%%%%%

% nguồn: Ok_TracNghiemLop11_Lan1
\begin{ex}%[1D7N3-1]%[ma:dac58507]
	Cho hàm số $f(x)=(2x-1)^{3}$.  Tính $f''(-1)$.
	\choice 
		{\True $f''(-1)=-72$} 
		{$f''(-1)=-27$} 
		 {$f''(-1)=-36$}
		{$f''(-1)=-18$} 
	 \loigiai{
		Ta có $f'(x)=6(2x-1)^{2}\Rightarrow  f''(x)=24(2x-1)\Rightarrow f''(-1)=-72$.
	 }
\end{ex}

% nguồn: Ok_Toan11_TracNghiem_Lan2
\begin{ex}%[1D7H3-1]%[ma:e24d65b2]
	 Đạo hàm cấp hai của hàm số $y=x^6-4x^3+2x+2\,022$ với $x\in\mathbb{R}$ là
	\choice 
		{$y''=30x^4-24x+2$} 
		{\True $y''=30x^4-24x$} 
		 {$y''=6x^5-12x^2+2$}
		 {$y''=6x^5-12x^2$}
	 \loigiai{
		Ta có  $y'=6x^5-12x^2+2$.\\
		Vậy  $y''=30x^4-24x$.
	 }
\end{ex}

% nguồn: Ok_Toan11_TracNghiem_Lan2
\begin{ex}%[1D7H3-1]%[ma:982debef]
	Cho  hàm số $f(x)=\dfrac{1}{2x-1}$. Tính $f''(-1)$.
	\choice 
		 {\True $-\dfrac{8}{27}$}
		 {$\dfrac{2}{9}$}
		 {$\dfrac{8}{27}$}
		 {$-\dfrac{4}{27}$}
	 \loigiai{
		Ta  có
		$f'(x)=\dfrac{-2}{(2x-1)^2}$  và $f''(x)=\dfrac{8}{(2x-1)^3}$.\\
		Vậy  $f''(-1)=-\dfrac{8}{27}$.
	} 
\end{ex}

% nguồn: Ok_Toan11_TracNghiem_Lan2
\begin{ex}%[1D7H3-1]%[ma:028975b8]
	Cho hàm số $y=\sin^2x$. Khi đó  $y''(x)$ bằng
	\choice 
		{$y''=\dfrac{1}{2}\cos 2x$} 
		{$P=2\sin  2x$}
		{\True  $y''=2\cos 2x$}
		{$y''=2\cos  x$}
	\loigiai{Ta có $y'=2\sin  x\cdot\cos x=\sin 2x$ và $y''=2\cos 2x$.
	 }
\end{ex}

% nguồn: Ok_Toan11_TracNghiem_Lan2
\begin{ex}%[1D7H3-1]%[ma:eba88751]
	Đạo hàm cấp hai của hàm  số $y=\cos^2x$ là
	\choice 
		 {\True $y''=-2\cos 2x$}
		{$y''=-2\sin  2x$}
		{$y''=2\cos  2x$}
		{$y''=2\sin  2x$}
	\loigiai{ 
		Ta có 
		$y'=2\cos x\cdot (-\sin x)=-\sin 2x$  và $y''=-2\cos 2x$.
	 }
\end{ex}

% nguồn: Ok_Toan11_TracNghiem_Lan2
\begin{ex}%[1D7H3-1]%[ma:b9a0e22f]
	Cho hàm số $y=x^5-3x^4+x+1$ với $x\in\mathbb{R}$. Đạo hàm $y''$ của hàm  số là
	\choice 
		 {$y''=5x^3-12x^2+1$}
		 {$y''=5x^4-12x^3$}
		 {$y''=20x^2-36x^3$}
		{\True  $y''=20x^3-36x^2$}
	 \loigiai{
		Ta  có $y'=5x^4-12x^3+1$ và $y''=20x^3-36x^2$.
	} 
\end{ex}

% nguồn: Ok_Toan11_TracNghiem_Lan2
\begin{ex}%[1D7H3-1]%[ma:496e51d2]
	Cho hàm số $y=\dfrac{x-2}{x+3}$.  Tính $y''$.
	 \choice
		 {$y''=\dfrac{-5}{(x+3)^3}$}
		{$y''=\dfrac{10}{(x+3)^2}$} 
		{\True  $y''=\dfrac{-10}{(x+3)^3}$}
		{$y''=\dfrac{5}{(x+3)^3}$} 
	\loigiai{ 
		 Ta có $y'=\dfrac{5}{(x+3)^2}$.\\
		Vậy  $y''=5\cdot\dfrac{-2(x+3)}{(x+3)^4}=\dfrac{-10}{(x+3)^3}$.
	} 
\end{ex}

% nguồn: Ok_Toan11_TracNghiem_Lan2
\begin{ex}%[1D7H3-1]%[ma:1f6d5a97]
	Tính đạo hàm cấp hai của hàm số $y=-3\cos x$ tại điểm $x_0=\dfrac{\pi}{2}$. 
	\choice 
		{$y''\left(\dfrac{\pi}{2}\right)=-3$} 
		 {$y''\left(\dfrac{\pi}{2}\right)=5$}
		{\True $y''\left(\dfrac{\pi}{2}\right)=0$} 
		{$y''\left(\dfrac{\pi}{2}\right)=3$} 
	\loigiai{ 
		Ta có 
		$y'=3\sin  x$ và $y''=3\cos x$.\\
		Vậy $y''\left(\dfrac{\pi}{2}\right)=0$. 
	} 
\end{ex}

% nguồn: Ok_Toan11_TracNghiem_Lan2
\begin{ex}%[1D7H3-1]%[ma:0af9747f]
	Đạo hàm  cấp hai của hàm số $y=\dfrac{3x+1}{x+2}$ là
	 \choice
		 {$y''=\dfrac{10}{(x+2)^2}$}
		{$y''=-\dfrac{5}{(x+2)^4}$} 
		{$y''=-\dfrac{5}{(x+2)^3}$} 
		 {\True $y''=-\dfrac{10}{(x+2)^3}$}
	 \loigiai{
		Ta có  $y'=\dfrac{5}{(x+2)^2}$ và $y''=-\dfrac{10}{(x+2)^3}$.
	} 
\end{ex}

% nguồn: Ok_Toan11_TracNghiem_Lan2
\begin{ex}%[1D7H3-1]%[ma:d1d83934]
	Cho hàm số $f(x)=x^3+2x$,  giá trị của $f''(1)$ bằng
	\choice 
		{\True $6$} 
		 {$8$}
		{$3$} 
		{$2$} 
	\loigiai{ 
		Ta có 
		$f'(x)=3x^2+2$  và $f''(x)=6x$.\\
		Vậy  $f''(1)=6$.
	 }
\end{ex}

% nguồn: Ok_Toan11_TracNghiem_Lan2
\begin{ex}%[1D7H3-1]%[ma:5ad6eeab]
	Cho hàm  số $y=-\dfrac{1}{x}$. Đạo hàm cấp hai của hàm số là
	 \choice
		{$y''=\dfrac{2}{x^3}$} 
		 {$y''=\dfrac{-2}{x^2}$}
		{\True $y''=\dfrac{-2}{x^3}$} 
		 {$y''=\dfrac{2}{x^2}$}
	 \loigiai{
		 Ta có $y'=\dfrac{1}{x^2}$ và $y''=-\dfrac{(x^2)'}{x^4}=-\dfrac{2x}{x^4}=-\dfrac{2}{x^3}$.
	} 
\end{ex}

% nguồn: Ok_Toan11_TracNghiem_Lan2
\begin{ex}%[1D7H3-1]%[ma:6be727c9]
	Cho hàm số  $f(x)=(3x-7)^5$. Tính $f''(2)$.
	 \choice
		 {$f''(2)=0$}
		 {$f''(2)=20$}
		{\True  $f''(2)=-180$}
		{$f''(2)=30$} 
	 \loigiai{Ta có $f'(x)=15(3x-7)^4$ và $f''(x)=180(3x-7)^3$.\\
		 Vậy $f''(2)=-180$.
	 }
\end{ex}

% nguồn: Ok_Toan11_TracNghiem_Lan2
\begin{ex}%[1D7H3-1]%[ma:4bc6be81]
	Cho $f(x)=x^3$.  Tính $f''(1)$.
	\choice 
		 {$f''(1)=3$}
		 {$f''(1)=2$}
		{\True $f''(1)=6$} 
		{$f''(1)=1$} 
	 \loigiai{
		Ta  có $f'(x)=3x^2$ và $f''(x)=3\cdot 2x=6x$.\\
		Vậy $f''(1)=6\cdot  1=6$.
	} 
\end{ex}

% nguồn: Ok_TracNghiemLop11_Lan1
\begin{ex}%[1D7H3-1]%[ma:4a223ded]
	Cho $f(x)=(2x-3)^5$. Khi  đó $f''(3)$ bằng
	\choice 
		{$810$} 
		{\True  $2160$}
		{$-2160$} 
		{$4320$} 
	\loigiai{Với $f(x)=(2x-3)^5$,  ta có
		 \allowdisplaybreaks
		\begin{eqnarray*}
			f'(x)&=&5\cdot(2x-3)^4\cdot(2x-3)'=5(2x-3)^4\cdot2=10(2x-5)^4,\\ 
			 f''(x)&=&10\cdot4\cdot(2x-3)^3\cdot(2x-3)'=80(2x-3)^3.
		\end{eqnarray*}
	Suy  ra $f''(3)=80(2\cdot3-3)^3=80\cdot3^3=2160$.}
\end{ex}

% nguồn: Ok_TracNghiemLop11_Lan1
\begin{ex}%[1D7H3-1]%[ma:be5593c7]
	Đạo hàm cấp hai của hàm  số $f(x)=\dfrac{4}{5}x^5-3x^2-x+4$ là
	 \choice
		{$16x^3-6x$} 
		{$4x^3-6$} 
		{\True  $16x^3-6$}
		{$4x^3-6x$} 
	\loigiai{$f'(x)=4x^4-6x-1,$ $f''(x)=16x^3-6.$ 
	 }
\end{ex}

% nguồn: Ok_TracNghiemLop11_Lan1
\begin{ex}%[1D7H3-1]%[ma:a96b4338]
	Đạo hàm cấp hai của hàm số  $ y=\dfrac{3x+1}{x+2}$ là
	 \choice
		 {$y''=\dfrac{10}{\left(x+2\right)^2}$}
		 {$y''=-\dfrac{5}{\left(x+2\right)^4}$}
		 {$y''=-\dfrac{5}{\left(x+2\right)^3}$}
		 {\True $y''=-\dfrac{10}{\left(x+2\right)^3}$}
	\loigiai{ 
		Ta  có $ y=3-\dfrac{5}{x+2}\Rightarrow{y}'=\dfrac{5}{\left(x+2\right)^2};{y}''=-\dfrac{10}{\left(x+2\right)^3}$\\
	} 
\end{ex}

% nguồn: Ok_TracNghiemLop11_Lan1
\begin{ex}%[1D7H3-1]%[ma:d5c4940e]
	Đạo hàm  cấp hai của hàm số $ y=x^6-4x^3+2x+2022$ với $ x\in\mathbb{R}$ là
	 \choice
		{$y''=30x^4-24x+2$} 
		 {$y''=30x^4-24x$}
		{\True $y''=6x^5-12x^2+2$} 
		 {$y''=6x^5-12x^2$}
	\loigiai{ 
	Ta  có $y'=6x^5-12x^2+2$. Suy ra $y''=30x^4-24x$.}
\end{ex}

% nguồn: Ok_TracNghiemLop11_Lan1
\begin{ex}%[1D7H3-1]%[ma:d6472d3c]
	Cho hàm số $f(x)=\sin^2x$, với mọi $x\in  \mathbb{R}$ ta có $f"(x)$ bằng
	\choice 
		{$f"(x)=2\cos x$} 
		 {$f"(x)=2\sin 2x$}
		 {$f"(x)=\cos 2x$}
		{\True $f"(x)=2\cos  2x$}
	\loigiai{ 
		Ta có: $f'(x)=2\sin x\cos x=\sin 2x$; $f"(x)=2\cos 2x$. 
	} 
\end{ex}


%%%%%%%%%%%%%%%% ĐÚNG/SAI %%%%%%%%%%%%%%%%

% nguồn: Ok_dungsai_lop11_Lan1
\begin{ex}%[1D7N3-1]%[ma:98496ae1]
 Cho hàm số $f(x)=x^3+2x$. Xét tính đúng, sai  của các mệnh đề sau
 \choiceTF 
  {\True Giá trị của $f''(1)$ bằng $6$} 
   {$f''(-5)>0$}
  {$3f''(2)-10=2$} 
  {\True $f''(x)>0\Leftrightarrow  x>0$}
 \loigiai{ 
  Ta  có $f'(x)=3x^2+2\Rightarrow f''(x)=6x$.
  \begin{itemize}
   \item $f''(1)=6$. 
   \item  $f''(-5)=-30<0$.
   \item $3f''(2)-10=3\cdot 6\cdot 2 - 10 =26\neq  2$.
   \item $f''(x)>0\Leftrightarrow 6x>0\Leftrightarrow x>0$. 
  \end{itemize}
 } 
\end{ex}

% nguồn: Ok_dungsai_lop11_Lan1
\begin{ex}%[1D7H3-1]%[ma:b6eaab1a]
 Xét tính  đúng, sai của các mệnh đề sau
  \choiceTF
   {\True Đạo hàm cấp hai của hàm số $y=\dfrac{2x^2-3x}{x-1}$ là $y''=\dfrac{-2}{(x-1)^3}$}
  {Cho hàm số $f(x)=x^3-3x^2+x+1$. Giá trị $f''(1)$  bằng $2$}
  {\True  Cho hàm số $f\left(x\right)=\cos x+3x^{2}$. Khi đó $f''\left(\dfrac{\pi }{2}\right)=6$}
   {Đạo hàm cấp hai của hàm số $y=-\dfrac{1}{x}$ là $y''=\dfrac{2}{x^3}$
  } 
 \loigiai{ 
  \begin{itemchoice}
   \itemch $y=2x-1-\dfrac{1}{x-1}\Rightarrow y'=2+\dfrac{1}{(x-1)^2}\Rightarrow y''=\dfrac{-2}{(x-1)^3}$. 
   \itemch Ta có  $f'(x)=3x^2-6x+1$, $f''(x)=6x-6$. Suy ra $f''(1)=0$.
   \itemch Ta có  $f'\left(x\right)=-\sin x+6x$ và $f''\left(x\right)=-\cos x+6$. Vậy $f''\left(\dfrac{\pi }{2}\right)=6$.
   \itemch Ta có $y'=\dfrac{1}{x^2}$ nên $y''=-\dfrac{(x^2)’}{x^4}=-\dfrac{2x}{x^4}=-\dfrac{2}{x^3}$. 
  \end{itemchoice}
  }
\end{ex}

% nguồn: Ok_dungsai_lop11_Lan1
\begin{ex}%[1D7H3-1]%[ma:e58bf365]
 Cho $y$ là hàm số của $x$ thỏa mãn phương trình $y=x^3$. Xét tính  đúng sai của các mệnh đề sau
  \choiceTF
  {Đạo hàm của hàm số $y$ là  $y'=3x^3$}
  {\True $y'(2)=12$} 
  {\True Đạo hàm cấp hai của hàm số $y$  là $y''=6x$}
   {$y''(0)=1$}
  \loigiai{
  \begin{itemchoice}
   \itemch  Sai vì $y'=3x^2$.
   \itemch Đúng vì $y'=3x^2$. Từ đó  $y'(2)=3\cdot 2^2=12$.
   \itemch  Đúng vì $y''=6x$.
   \itemch Sai vì $y''(0)=6\cdot  0=0$.
  \end{itemchoice}
 } 
\end{ex}

% nguồn: Ok_dungsai_lop11_Lan1
\begin{ex}%[1D7H3-1]%[ma:8d596c57]
 Cho $y$ là hàm số  của $x$ thỏa mãn phương trình $y=\sin x$. Xét tính đúng sai của các mệnh đề sau
 \choiceTF 
   {\True Đạo hàm của hàm số $y$ là $y'=\cos x$}
   {\True $y'\left(\dfrac{\pi}{6}\right)=\dfrac{\sqrt{3}}{2}$}
  {Đạo hàm cấp hai của hàm  số $y$ là $y''=-\sin^2x$}
   {$y''\left(\dfrac{\pi}{2}\right)=1$}
 \loigiai{ 
  \begin{itemchoice}
   \itemch Đúng  vì $y'=\cos x$.
   \itemch Đúng vì $y'=\cos x$. Từ đó  $y'\left(\frac{\pi}{6}\right)=\cos\left(\dfrac{\pi}{6}\right)=\dfrac{\sqrt{3}}{2}$.
    \itemch Sai vì $y'=\cos x$; $y''=-\sin x$.
   \itemch Sai vì Ta có $y''\left(\dfrac{\pi}{2}\right)=-\sin \dfrac{\pi}{2}=-1$. 
 \end{itemchoice}}
\end{ex}

% nguồn: Ok_dungsai_lop11_Lan1
\begin{ex}%[1D7H3-1]%[ma:8d081ab9]
 Cho hàm số $y=f(x)=2\sin x+3\cos x$. Xét tính đúng  sai của các mệnh đề sau
  \choiceTF
  {\True  $f'(x)=2\cos x-3\sin x, \forall x\in\mathbb{R}$}
  {$f'(0)=-2$} 
  {\True $f''(x)=-2\sin x-3\cos x,  \forall x\in\mathbb{R}$}
   {$f''\left(\dfrac{\pi}{2}\right)=2$}
  \loigiai{
  \begin{itemchoice}
   \itemch Đúng  vì $f'(x)=2(\sin x)'+3(\cos x)'=2\cos x+3\cdot (-\sin x)=2\cos x+-3\sin x$
   \itemch  Sai vì $f'(0)=2\cos 0-3\sin 0=2$.
   \itemch Đúng vì $f''(x)=2(\cos x)'-3(\sin x)'=-2\sin x-\cos x,  \forall x\in\mathbb{R}$.
   \itemch  Sai vì $f''\left(\dfrac{\pi}{2}\right)=-2\sin\left(\dfrac{\pi}{2}\right)-\cos\left(\dfrac{\pi}{2}\right)=-2$.
  \end{itemchoice}
 } 
\end{ex}

% nguồn: Ok_dungsai_lop11_Lan1
\begin{ex}%[1D7H3-1]%[ma:644d80e8]
 Cho hàm số $y=\ln  x$ với $x>0$. Xét tính đúng sai của các mệnh đề sau
 \choiceTF 
  {\True Đạo hàm của  hàm số $y$ là $y'=\dfrac{1}{x}$}
   {Đạo hàm cấp hai của hàm số $y$ là $y''=\dfrac{-1}{x^4}$ với $x>0$}
   {\True $y'(1)=1$}
  {$y''(2)=\dfrac{-1}{4}$} 
 \loigiai{ 
  \begin{itemchoice}
   \itemch Đúng vì $y=\ln x$; $y'=\dfrac{1}{x}$ với $x>0$. 
   \itemch Sai vì $y=\ln x$; $y'=\dfrac{1}{x}$;  $y''=\dfrac{-1}{x^2}$ với $x>0$.
   \itemch Đúng  vì $y'(1)=1$.
   \itemch Đúng.  $y''(2)=\dfrac{-1}{4}$.
  \end{itemchoice}
  }
\end{ex}

% nguồn: Ok_dungsai_lop11_Lan1
\begin{ex}%[1D7H3-1]%[ma:d796a2e3]
 Cho hàm số $y=2\sin 2x-\cos\left(x-\dfrac{\pi}{3}\right)$. Xét tính đúng sai  của các mệnh đề sau
  \choiceTF
   {$y'\left(\dfrac{\pi}{3} \right)=2$}
  {\True $y''=-8\cdot\sin 2x+\cos\left(x-\dfrac{\pi}{3}\right)$}{$y''(0)=1$} 
  {$y'\left(\dfrac{\pi}{3}\right)-2y''(0)=-2$} 
  \loigiai{
  Ta  có $y'=4\cdot\cos 2x+\sin\left(x-\dfrac{\pi}{3}\right)$.
  \begin{itemchoice}
   \itemch Sai vì $y'\left(\dfrac{\pi}{3}\right)=4\cdot\cos  \dfrac{2\pi}{3}+\sin 0=-2$.
   \itemch Đúng. 
   \itemch Sai vì 
   $y''(0)=\dfrac{1}{2}  $.
    \itemch Sai vì $y'\left(\dfrac{\pi}{3}\right)-2y''(0)=-2-2\cdot\dfrac{1}{2}=-3$.
  \end{itemchoice}
  }
\end{ex}

% nguồn: Ok_dungsai_lop11_Lan1
\begin{ex}%[1D7H3-1]%[ma:c0b42162]
 Biết $y$ là hàm số của $x$ thỏa mãn phương trình $xy=1+\ln y$. Xét tính đúng  sai của các mệnh đề sau
  \choiceTF
  {Đạo hàm của hàm số  $y$ là $y'=\dfrac{y^2}{1+xy}$}
   {\True $y'(0)=\dfrac{1}{e^2}$}
  {\True  Đạo hàm cấp hai của hàm số $y$ là $y''=\dfrac{2yy'(1-xy)+y^2\left(y+xy'\right)}{(1-xy)^2}$}
  {\True  $y''(0)=\dfrac{3}{e^3}$}
 \loigiai{ 
  \begin{itemchoice}
   \itemch Sai. Ta có $xy=1+\ln  y\Rightarrow y+xy'=\dfrac{y'}{y}\Rightarrow y'=\dfrac{y^2}{1-xy}$.
   \itemch Đúng. Thay $x=0$ vào  phương trình $xy=1+\ln y$ ta được $1+\ln y=0\Rightarrow y=\dfrac{1}{e}$.\\
    Vậy $y'(0)=\dfrac{\left(\frac{1}{e}\right)^2}{1-0\cdot\frac{1}{e}}=\dfrac{1}{e^2}$.
   \itemch Đúng. Ta  có $y''=\left(\dfrac{y^2}{1-xy}\right)'=\dfrac{2yy'(1-xy)+y^2\left(y+xy'\right)}{(1-xy)^2}$.
   \itemch Đúng. Thay $y=\dfrac{1}{e}$; $x=0$; $y'=\dfrac{1}{e^2}$ vào  $y''$ ta được $y''(0)=\dfrac{2\cdot\frac{1}{e}\cdot\frac{1}{e^2}+\frac{1}{e^2}\cdot\frac{1}{e}}{1}=\dfrac{3}{e^3}$.
  \end{itemchoice}
  }
\end{ex}

% nguồn: Ok_dungsai_lop11_Lan2
\begin{ex}%[1D7H3-1]%[ma:871cb9e7]
	 Cho hàm số $y=2\sin 2x-\cos\left(x-\dfrac{\pi}{3}\right)$.
	\choiceTF 
		{$y'\left(\dfrac{\pi}{3}\right)=2$} 
		 {\True $y''=-8\cdot\sin 2x+\cos\left(x-\dfrac{\pi}{3}\right)$}
		{$y''(0)=1$} 
		{$y'\left(\dfrac{\pi}{3}\right)-2y''(0)=-2$} 
	 \loigiai{
		\begin{itemchoice}
			\itemch Ta có  $y'=4\cos 2x+\sin\left(x-\dfrac{\pi}{3}\right)\Rightarrow y'\left(\dfrac{\pi}{3}\right)=4\cdot\cos\dfrac{2\pi}{3}+\sin 0=-2$.
			\itemch Ta có $y''=-8\cdot\sin  2x+\cos\left(x-\dfrac{\pi}{3}\right)$.
			\itemch Ta có  $y''=-8\cdot\sin 2x+\cos\left(x-\dfrac{\pi}{3}\right)\Rightarrow y''(0)=\dfrac{1}{2}$.
			\itemch Ta có $y'\left(\dfrac{\pi}{3}\right)-2y''(0)=-2-2\cdot\dfrac{1}{2}=-3$. 
		\end{itemchoice}
	 }
\end{ex}

% nguồn: Ok_dungsai_lop11_Lan2
\begin{ex}%[1D7H3-1]%[ma:8bdc33e8]
	Cho hàm  số $f(x)=\dfrac{1}{12}x^4-\dfrac{5}{6}x^3+3x^2+3x$.
	\choiceTF 
		{$f''(1)=0$} 
		{\True  $f''(3)=0$}
		{\True $f''(x) < 0\Leftrightarrow 2< x < 3$} 
		{Đồ thị hàm số $y=f''(x)$ có trục  đối xứng là đường thẳng $x=5$}
	 \loigiai{
		Ta  có $f'(x)=\dfrac{1}{3}x^3-\dfrac{5}{2}x^2+6x+3$ và $f''(x)=x^2-5x+6$.
		\begin{itemchoice}
			\itemch Ta có  $f''(1)=1^2-5\cdot 1+6=2$.
			\itemch Ta có $f''(3)=0$. 
			\itemch Ta  có $f''(x)=x^2-5x+6<0\Leftrightarrow 2<x<3$.
			\itemch Đồ thị hàm số  $y=f''(x)$ là một parabol có trục đối xứng là $x=-\dfrac{b}{2a}=\dfrac{5}{2}$.
		\end{itemchoice}
	} 
\end{ex}


%%%%%%%%%%%%%%%% TRẢ LỜI NGẮN %%%%%%%%%%%%%%%%

% nguồn: Ok_TraLoiNgan_Lop11_Lan1
\begin{ex}%[1D7H3-1]%[ma:9de32aee]
	 Đạo hàm cấp hai của hàm số $f(x)=\dfrac{1}{5}x^{5}-3 x^{2}-x+4$ là $f''(x)=mx^3+n$. Tính $m+n^2$.
	\shortans{40} 
	 \loigiai{
		Ta có $ y'=x^4-6x-1\Rightarrow y''=4x^3-6 \Rightarrow m=4, n=-6$.\\ Vậy $m+n^2=4+(-6)^2=40$. 
	 }
\end{ex}

% nguồn: Ok_TraLoiNgan_Lop11_Lan1
\begin{ex}%[1D7H3-1]%[ma:81addbac]
	 Cho hàm số $f(x)=\sin^2x-x^2+1$. Tính $f^{\prime \prime}\left(\dfrac{\pi}{2}\right)$.
	 \shortans{-4}
	 \loigiai{
		$f^\prime (x)=2\sin x \cdot  \cos x-2x=\sin2x -2x$.\\
		$f^{\prime \prime}(x)=2\cos2x-2$.  Vậy
		$f^{\prime \prime}\left(\dfrac{\pi}{2}\right)=2\cos  \pi -2=-4$.
	 }
\end{ex}

% nguồn: Ok_TraLoiNgan_Lop11_Lan1
\begin{ex}%[1D7H3-1]%[ma:69bb550a]
	Có bao nhiêu cặp số nguyên dương $(a; b)$ thỏa mãn $\log_{a}b+6 \log_{b} a=5$  và $a \in[2; 2005]$, $b \in[2;2005]$?
	 \shortans{54}
	 \loigiai{
		Với các điều kiện trên, ta có 
		\begin{eqnarray*}
			&&\log_{a}b+6 \log_{b} a=5 
			 \\
			&\Leftrightarrow& \log_a  b +6 \dfrac{1}{\log_a b}=5
			\\ 
			&\Leftrightarrow&  \log^2_a b -5\log_a b+6=0
			 \\
			&\Leftrightarrow& \hoac{& \log_a b  =2 \\ & \log_a b =3}
			 \\
			&\Leftrightarrow& \hoac{&  b=a^2 \\ & b=a^3.}
		\end{eqnarray*}
		Khi $b=a^2 \Rightarrow 2\le a^2 \le 2005 \Rightarrow \sqrt  2 \le a \le \sqrt{2005}$.
		\\ 
		Do $a$ nguyên nên $2\le a \le 44$, vậy có $44-2+1=43$ số $a$ thoả mãn nên có $43$  cặp số $a$, $b$ thoả mãn yêu cầu bài toán.
		 \\
		Khi $b=a^3 \Rightarrow 2\le a^3 \le 2005 \Rightarrow \sqrt[3] 2 \le a \le  \sqrt[3]{2005}$.
		 \\
		Do $a$ nguyên nên $2\le a \le 12$, vậy có $12-2+1=11$ số $a$ thoả mãn nên có $11$ cặp số $a$, $b$  thoả mãn yêu cầu bài toán.
		\\ 
		Vậy có tổng cộng $43+11=54$ cặp  số thoả mãn.
	 }
\end{ex}

% nguồn: Ok_TraLoiNgan_Lop11_Lan2
\begin{ex}%[1D7H3-1]%[ma:c5c3b22a]
	Tính đạo hàm cấp hai của hàm số $y=x\cdot 2^{2x}$  tại $x_0=0$ (kết quả làm tròn đến hàng phần trăm).
	 \shortans{2,77}
	\loigiai{ 
		Ta có $ y'=2^{2x}+x\cdot 2\cdot 2^{2x}\ln{2} =(2\ln 2  \cdot x+1)\cdot 2^{2x}$.\\
		Khi đó $ y''=\left((2x+1)\cdot 2^{2x}\right)'=2\ln  2\cdot 2^{2x}+(2x+1)\cdot 2 \cdot 2^{2x}\cdot \ln 2 =(2\ln 2 +4x\ln 2 +2\ln 2)\cdot 2^{2x}=(4x\ln2+4\ln 2)\cdot 2^{2x}$.\\
		 Vậy $ y''(0)=4\ln 2 \approx 2{,}77 $.
	 }
\end{ex}

% nguồn: Ok_TraLoiNgan_Lop11_Lan2
\begin{ex}%[1D7H3-1]%[ma:9c3d3390]
	Tính đạo hàm cấp hai của  hàm số $y=x\sin x$ tại $x_0=0$.
	\shortans{2} 
	 \loigiai{
		Ta có 
		\begin{itemize}
			\item $y'=(x\sin x)'=(x)'\cdot\sin x+(\sin  x)'\cdot x=\sin x+x\cos x$,
			\item  $y''=(\sin x+x\cos x)'=(\sin x)'+(x\cos x)'=\cos x+(x)'\cdot\cos x+(\cos x)'\cdot x=2\cos x-x\sin x$.
		\end{itemize}
		Vậy $y''(0)=2\cos  0-0\cdot\sin 0=2$.
	 }
\end{ex}

% nguồn: Ok_TraLoiNgan_Lop11_Lan2
\begin{ex}%[1D7V3-1]%[ma:95d782f0]
	Biết hàm số $y=f(x)-f\left(3x\right)$ liên tục trên và  có đạo hàm cấp hai với mọi $ x \in \mathbb{R} $, biết $y''(3)=5$, $y''(1)=10$. Tính đạo hàm cấp hai của hàm số $g(x)=f(x)-f\left(9x\right)$ tại $x=1$.
	\shortans{55} 
	\loigiai{ 
		Đặt  $h(x)=f(x)-f\left(3x\right)$; $g(x)=f(x)-f\left(9x\right)$.\\
		 Ta có $h'(x)=f'(x)-3f'\left(3x\right)\Rightarrow h''(x)=f''(x)-9f''\left(3x\right)$;\\
		$g'(x)=f'(x)-9f'\left(9x\right)\Rightarrow g''(x)=f''(x)-81f''\left(9x\right)\Rightarrow g''(1)=f''(1)-81f''(9)$ .\\ 
		Ta có $\heva{&h''(3)=f''(3)-9f''(9)=5 \Leftrightarrow  9f''(3)-81f''(9)=45 \qquad (1) \\
		 &h''(1)=f''(1)-9f''(3) =10.\qquad (2)}$\\
	Cộng từng vế $ (1) $ và $  (2) $ suy ra $g''(1)=f''(1)-81f''(9)=55$.}
\end{ex}


%%%%%%%%%%%%%%%% TỰ LUẬN %%%%%%%%%%%%%%%%

% nguồn: Ok_TuanLuan_Toan11_Lan2
\begin{ex}%[1D7H3-1]%[ma:1628c02e]
	Tìm đạo hàm cấp hai của các hàm số  sau:
	\begin{enumEX}{4}
		\item $y=\dfrac{1}{1-x}$, 
		\item  $y=\dfrac{1}{\sqrt{1-x}}$,
		\item $y=\tan  x$,
		 \item $y=\cos^2x$.
	\end{enumEX}
	 \loigiai{
		\begin{enumerate}
			 \item Ta có $y'=\dfrac{1}{(1-x)^2}$ và $y''=\dfrac{2}{(1-x)^3}$.
			\item Ta có $y'=\dfrac{1}{2(1-x)\sqrt{1-x}}$ và $y''=\dfrac{3}{4(1-x)^2\sqrt{1-x}}$. 
			\item Ta có $y'=\dfrac{1}{\cos^2x}$  và $y''=\dfrac{2\sin x}{\cos^3x}$.
			\item Ta có $y=\dfrac{1}{2}+\dfrac{1}{2}\cos 2x$. Khi đó $y'=-\sin 2x$ và $y''=-2\cos  2x$.
		\end{enumerate}
	 }
\end{ex}

% nguồn: Ok_TuanLuan_Toan11_Lan2
\begin{ex}%[1D7H3-1]%[ma:ee0fae64]
	Tính đạo hàm cấp  hai của hàm số $y=\dfrac{x+1}{x-2}$.
	 \loigiai{
		Ta  có $y=\dfrac{x+1}{x-2}=1+\dfrac{3}{x-2}\Rightarrow y'=-\dfrac{3}{(x-2)^2}$ và $y''=\dfrac{6}{(x-2)^3}$.
	} 
\end{ex}

% nguồn: Ok_TuanLuan_Toan11_Lan2
\begin{ex}%[1D7H3-1]%[ma:a69cb4d5]
	Cho $f(x)=(x+10)^6$. Tính  $f''(2)$.
	 \loigiai{
		 Ta có $f'(x)=6(x+10)^5$.\\
		Suy ra $f''(x)=6\cdot 5\cdot (x+10)^4=30(x+10)^4$.\\ 
		 Vậy $f''(2)=30\cdot (2+10)^4=622\,080$.
	} 
\end{ex}

% nguồn: Ok_TuanLuan_Toan11_Lan2
\begin{ex}%[1D7H3-1]%[ma:171a4ab8]
	Cho  $f(x)=\sin 3x$. Tính $f''\left(-\dfrac{\pi}{2}\right)$, $f''(0)$, $f''\left(\dfrac{\pi}{18}\right)$.
	 \loigiai{
		 Ta có $f'(x)=3\cos 3x$ và $f''(x)=-9\sin 3x$.\\
		Khi đó 
		\begin{itemize}
			\item $f''\left(-\dfrac{\pi}{2}\right)=-9\sin\left(-\dfrac{3\pi}{2}\right)=-9$; 
			 \item $f''(0)=-9\sin (0)=0$;
			\item  $f''\left(\dfrac{\pi}{18}\right)=-9\sin\left(\dfrac{\pi}{6}\right)=\dfrac{-9}{2}$.
		\end{itemize}
	 }
\end{ex}

% nguồn: Ok_TuanLuan_Toan11_Lan2
\begin{ex}%[1D7H3-1]%[ma:3c6f2851]
	Tìm đạo hàm cấp hai  của các hàm số sau:
	\begin{enumEX}{2}
		\item $y=\sin 5x\cos 2x$, 
		\item  $y=\dfrac{2x+1}{x^2+x-2}$,
		\item $y=\dfrac{x}{x^2-1}$, 
		 \item $y=\dfrac{x-1}{x+2}$,
		\item  $y=x^2\sin x$,
		\item $y=x\sqrt{1+x^2}$, 
		\item $y=\left(1-x^2\right)\cos x$, 
		\item $y=\sqrt{x}$, 
		\item $y=\sin  x\sin 2x\sin 3x$,
		\item  $y=\dfrac{x^2}{1-x}$
		\item  $y=x\cos 2x$,
		\item  $y=\dfrac{1}{\sqrt{x}}$.
	\end{enumEX}
	 \loigiai{
		\begin{enumerate}
			 \item
			Ta có $y=\sin 5x\cos 2x=\dfrac{1}{2}(\sin 7x+\sin 3x)$.\\ 
			Khi  đó $y'=\dfrac{1}{2}(7\cos 7x+3\cos 3x)$ và $y''=-\dfrac{1}{2}(49\sin 7x+9\sin 3x)$.
			\item  Ta có $y=\dfrac{2x+1}{x^2+x-2}=\dfrac{1}{x-1}+\dfrac{1}{x+2}$.\\
			Khi đó $y'=-\dfrac{1}{(x-1)^2}-\dfrac{1}{(x+2)^2}$ và $y''=\dfrac{2}{(x-1)^3}+\dfrac{2}{(x+2)^3}$. 
			\item  Ta có $y=\dfrac{x}{x^2-1}=\dfrac{1}{2}\left(\dfrac{1}{x+1}+\dfrac{1}{x-1}\right)$.\\
			Khi đó  $y'=\dfrac{1}{2}\left(\dfrac{-1}{(x+1)^2}+\dfrac{-1}{(x-1)^2}\right)$ và $y''=\dfrac{1}{(x+1)^3}+\dfrac{1}{(x-1)^3}$.
			 \item Ta có $y=\dfrac{x-1}{x+2}=1-\dfrac{3}{x+2}$.\\
			Khi đó $y'=\dfrac{3}{(x+2)^2}$  và $y''=-\dfrac{6}{(x+2)^3}$.
			\item Ta có $y'=2x\cdot\sin x+x^2\cdot\cos x$ và \begin{align*}
			y''&=\left(2x\cdot\sin x\right)'+\left(x^2\cdot\cos x\right)'=2\sin x+2x\cos x+2x\cos x-x^2\sin  x\\
			&=\left(2-x^2\right)\sin x+4x\cos x. 
		\end{align*}
		\item Ta có $y'=\dfrac{2x^2+1}{\sqrt{1+x^2}}$ và \begin{align*}
		y''&=\left(\dfrac{2x^2+1}{\sqrt{1+x^2}}\right)'=\dfrac{4x\cdot\sqrt{1+x^2}-\left(2x^2+1\right)\cdot\dfrac{x}{\sqrt{1+x^2}}}{1+x^2}\\ 
		&=\dfrac{4x\left(1+x^2\right)-\left(2x^2+1\right)x}{\left(1+x^2\right)\sqrt{1+x^2}}=\dfrac{2x^3+3x}{\left(1+x^2\right)\sqrt{1+x^2}}. 
	\end{align*}
	 \item Ta có $y'=-2x\cos x-\left(1-x^2\right)\sin x$ và
	\begin{align*}
		y''&=\left(-2x\cos x\right)'-\left[\left(1-x^2\right)\sin  x\right]'=-2\cos x+2x\sin x+2x\sin x-\left(1-x^2\right)\cos x\\
		&=\left(x^2-3\right)\cos x+4x\sin x. 
	\end{align*}
	\item Ta có $y'=\dfrac{1}{2\sqrt{x}}$  và $y''=-\dfrac{1}{4x\sqrt{x}}$.
	\item Ta có $y=\sin x\sin 2x\sin 3x=\dfrac{1}{4}\sin 2x+\dfrac{1}{4}\sin 4x-\dfrac{1}{4}\sin  6x$.\\
	Khi đó $y'=\dfrac{1}{2}\cos 2x+\cos 4x-\dfrac{3}{2}\cos  6x$ và $y''=-\sin 2x-4\sin 4x+9\sin 6x$.
	\item Ta có  $y=\dfrac{x^2}{1-x}=-x-1+\dfrac{1}{1-x}$.\\
	 Khi đó $y'=-1+\dfrac{1}{(1-x)^2}$ và $y''=\dfrac{2}{(1-x)^3}$.
	\item Ta có  $y'=\cos 2x-2x\sin 2x$ và $y'=-4\sin 2x-4x\cos 2x$.
	\item  Ta có $y'=-\dfrac{1}{2x\sqrt{x}}$ và $y''=\dfrac{3}{4\sqrt{x^5}}$.
	\end{enumerate}
 }
\end{ex}

% nguồn: Ok_TuanLuan_Toan11_Lan2
\begin{ex}%[1D7H3-1]%[ma:b09c41cc]
	Tính đạo hàm cấp hai của hàm số $f(x)=x^3-x^2+1$ tại  điểm $x=2$.
	\loigiai{ 
		Ta  có $f'(x)=3x^2-2x$.\\
		 Suy ra $f''(x)=6x-2$.\\
		 Vậy $f''(2)=10$.
	} 
\end{ex}
